\DocumentMetadata{
  pdfversion=2.0,
  pdfstandard=ua-2,
  lang=en-US,
  tagging=on,
  tagging-setup={math/setup=mathml-SE,tabsorder=structure}
}
\documentclass[11pt,letterpaper]{article}
\usepackage[margin=0.68in,includefoot]{geometry}
\usepackage{newcomputermodern}
\usepackage{amsmath,amscd,mathtools}
\providecommand{\square}{\mdlgwhtsquare}
\usepackage[hidelinks]{hyperref}
\hypersetup{
  pdftitle={Analysis PhD Exam, May 2014},
  pdfauthor={University of Florida Department of Mathematics},
  pdfsubject={Graduate mathematics examination},
  pdfdisplaydoctitle=true,
  bookmarksopen=true
}
\setcounter{secnumdepth}{0}
\setlength{\parindent}{0pt}
\setlength{\parskip}{0.28em}
\setlength{\emergencystretch}{3em}
\setlength{\leftmargini}{2.15em}
\setlength{\leftmarginii}{2.65em}
\setlength{\leftmarginiii}{3em}
\renewcommand{\labelenumi}{\textbf{\theenumi.}}
\pagestyle{empty}
\raggedbottom
\allowdisplaybreaks
\newenvironment{examproblems}[1]{%
  \begin{enumerate}%
  \setcounter{enumi}{#1}%
  \setlength{\topsep}{0.35em}%
  \setlength{\partopsep}{0pt}%
  \setlength{\itemsep}{0.48em}%
  \setlength{\parsep}{0.18em}%
}{\end{enumerate}}
\newenvironment{examsubparts}{%
  \begin{enumerate}%
  \setlength{\topsep}{0.18em}%
  \setlength{\partopsep}{0pt}%
  \setlength{\itemsep}{0.16em}%
  \setlength{\parsep}{0.1em}%
}{\end{enumerate}}
\newenvironment{examinstructions}{%
  \par\begingroup\itshape
}{%
  \par\endgroup\vspace{0.18em}
}
\makeatletter
\renewcommand\section{\@startsection{section}{1}{0pt}%
  {-1.25ex plus -.25ex minus -.1ex}%
  {0.55ex plus .1ex}%
  {\normalfont\large\bfseries}}
\renewcommand{\@maketitle}{%
  \newpage\null\vskip -1.2em%
  {\footnotesize\bfseries
    \href{https://www.ufl.edu/}{University of Florida}, \href{https://math.ufl.edu/}{Department of Mathematics}\par}%
  \vskip 0.18em%
  \tagstructbegin{tag=Title}%
    {\LARGE\bfseries\href{https://gma.math.ufl.edu/exams/analysis-phd/}{Analysis PhD Exam}, May 2014\par}%
  \tagstructend%
  \vskip 0.35em%
  \tagmcbegin{artifact}%
    \hrule height 1pt%
  \tagmcend%
  \vskip 0.55em%
}
\makeatother
\title{Analysis PhD Exam, May 2014}
\author{}
\date{}
\begin{document}
\maketitle
\thispagestyle{empty}
\begin{examinstructions}
You must answer SIX questions. Write solutions in a neat and logical fashion, giving complete reasons for all steps.
\end{examinstructions}
\begin{examproblems}{0}
\item
Let \((X,\mathcal{M},\mu)\) be a measure space and \(f\in L^1(\mu)\). Prove that for every \(\epsilon>0\), there exists a \(\delta>0\) such that if \(\mu(E)<\delta\), then \(\int_E |f|\,d\mu<\epsilon\).
\item
State the Hahn decomposition theorem for signed measures. Prove that if~\(\rho\) is a signed measure on a measurable space \((X,\mathcal{M})\), then there exist unique positive measures \(\rho_+,\rho_-\) such that \(\rho_+\perp\rho_-\) and \(\rho=\rho_+-\rho_-\).
\item
State the Fubini–Tonelli theorem and give a sketch of its proof.
\item
Let \(C^1[0,1]\) denote the set of all continuous (real-valued) functions~\(f\) on \([0,1]\) such that~\(f\) is differentiable in \((0,1)\) and \(f'\) extends continuously to \([0,1]\). Prove that \(C^1[0,1]\) is a Banach space under the norm 
\[
\|f\|=\|f\|_\infty+\|f'\|_\infty.
\]
\item
Let~\(\mathcal{X}\) be a normed vector space (over~\(\mathbb{C}\)) and \(\mathcal{M}\subset\mathcal{X}\) a closed subspace. Prove that if \(x\in\mathcal{X}\setminus\mathcal{M}\), then \(\mathcal{M}+\mathbb{C}x\) is closed.
\item
This problem has two parts.
\begin{examsubparts}
\item[\textnormal{a)}]
State the Closed Graph Theorem and the Banach Isomorphism Theorem.
\item[\textnormal{b)}]
Assuming the Banach Isomorphism Theorem (or otherwise), prove the Closed Graph Theorem.
\end{examsubparts}
\item
Let~\(\mathcal{X}\) be a Banach space and~\(\mathcal{Y}\) a normed vector space. Suppose that \(T_n\) is a sequence of bounded linear operators from~\(\mathcal{X}\) to~\(\mathcal{Y}\). Prove that if \(\lim T_nx\) exists for each \(x\in\mathcal{X}\), then the mapping \(Tx=\lim T_nx\) defines a bounded linear operator from~\(\mathcal{X}\) to~\(\mathcal{Y}\).
\item
Let \((X,\mathcal{M},\mu)\) be a~\(\sigma\)-finite measure space. Prove that the simple functions that belong to \(L^2(\mu)\) are dense in \(L^2(\mu)\).
\end{examproblems}
\end{document}
